% =====================================================================
%  YUYOS DE LA QUEBRADA DE CAFAYATE — DATOS QUE COMPLETA LA EDITORIAL
%
%  Este es el único archivo que hace falta tocar para cerrar la edición.
%  Lo leen cafayate.tex (página de créditos y colofón) y cubierta.tex
%  (lomo, sello y código de barras).
%
%  Mientras un dato esté vacío, el PDF lo muestra en un recuadro
%  amarillo y la compilación avisa cuántos faltan. Se completa
%  escribiendo el dato entre las llaves, sin tocar nada más.
%  Qué es cada dato y quién lo provee: PENDIENTES-editorial.md.
% =====================================================================

% ------------------------ Página de créditos (p. iv) -----------------
\def\datoSello{}           % pie editorial. P. ej.: Editorial Juana Manuela, Salta
\def\datoISBN{}            % con guiones. P. ej.: 978-987-XXXX-XX-X

% La autora confirmó por escrito la licencia CC BY-NC-SA 4.0 de la p. iv:
% cambiar \licenciaconfirmadafalse por \licenciaconfirmadatrue.
% (Si eligió otra licencia, hay que cambiar además el texto de la p. iv.)
\newif\iflicenciaconfirmada
\licenciaconfirmadafalse

% ------------------------------ Colofón (p. 45) ----------------------
\def\datoFechaImpresion{}  % P. ej.: noviembre de 2026
\def\datoImprenta{}        % nombre y ciudad. P. ej.: Imprenta X, Salta
\def\datoPapel{}           % interior. P. ej.: obra de 90 g/m²
\def\datoTapas{}           % P. ej.: cartulina ilustración de 300 g/m², laminado mate
\def\datoTirada{}          % solo el número. P. ej.: 500

% ------------------------------- Cubierta ----------------------------
% Espesor de UNA hoja (dos páginas) del papel del interior, en mm y con
% punto decimal. Lo da la imprenta. P. ej.: 0.11
% Mientras esté vacío, el lomo se calcula con 0.11 mm, provisorio.
\def\datoEspesorHoja{}

% Código de barras del ISBN: lo genera la editorial. Guardarlo en la
% carpeta del libro como codigo-barras-isbn.pdf (o .png), y la cubierta
% lo ubica solo en su lugar de la contratapa (50 x 30 mm).
